% ЗГЕНЕРОВАНО scripts/build_format_v2_chapter.py — не редагувати вручну.
% Джерело: editorial/cards/8.*.json, format-v2/rozdil-8.json, paraleli.json.

\naukchapter{8}{4}{Глава 8. Полон і визволення — в розумі}
\naukintro{Восьма глава — чотири вірші, дві пари. Перша пара (8.1 і 8.2) побудована як дослівне заперечення: ті самі шість дієслів, у другому вірші кожне з часткою na. Друга пара міняє мірило: спершу привʼязаність до поглядів (8.3), далі саме слово я (8.4). Визволення тут не додає нічого — воно знімає рівно те, що було названо полоном.}
\vsec{8.1}
\vdeva{तदा बन्धो यदा चित्तं किञ्चिद् वाञ्छति शोचति । किञ्चिन् मुञ्चति गृह्णाति किञ्चिद् धृष्यति कुप्यति ॥}
\viast{tadā bandho yadā cittaṃ kiñcid vāñchati śocati | kiñcin muñcati gṛhṇāti kiñcid dhṛṣyati kupyati ||}
\vtranslit{тада бандго яда чіттам кіньчід ваньччаті шочаті | кіньчін муньчаті ґріхнаті кіньчід дгрішяті куп'яті ||}
\vperek{Полон — тоді, коли розум чогось бажає, сумує, щось відпускає або бере, чомусь радіє чи гнівається.}
\vcom{Построкове зіставлення: редакторська чернетка, не рецензовано.}
\vcom{Рядок а}
\vdeva{तदा बन्धो यदा चित्तं किञ्चिद् वाञ्छति शोचति}
\viast{tadā bandho yadā cittaṃ kiñcid vāñchati śocati}
\vtranslit{тада бандго яда чіттам кіньчід ваньччаті шочаті}
\vperek{Полон — тоді, коли розум чогось бажає, сумує,}
\vcom{Кореляція tadā ... yadā пояснює, коли постає полон. У першому піврядку розум бажає й сумує; це два з шести дієслів переліку. Переклад залишає речення відкритим до наступного рядка.}
\vcom{Рядок б}
\vdeva{किञ्चिन् मुञ्चति गृह्णाति किञ्चिद् धृष्यति कुप्यति}
\viast{kiñcin muñcati gṛhṇāti kiñcid dhṛṣyati kupyati}
\vtranslit{кіньчін муньчаті ґріхнаті кіньчід дгрішяті куп'яті}
\vperek{щось відпускає або бере, чомусь радіє чи гнівається.}
\vcom{Другий піврядок завершує перелік: розум відпускає, бере, радіє та гнівається. Форма dhṛṣyati в джерельній картці не має остаточного лексикографічного розв'язання; відповідник «радіє» залишено без самовільного виправлення санскриту.}
\vcom{Визначення полону тут не метафізичне, а **описове**: шість дієслів, і кожне — звичайний рух розуму. Зверніть увагу, що «бере» і «відпускає» стоять в одному ряду з гнівом: вибір сам по собі є ознакою полону. Перелік не оцінює, а перелічує — і саме тому діє.}
\vpar{Українська паралель (наша інтерпретація)}{Григорій Сковорода, «Сад божественних пісень».}{Наша паралель: у Сковороди буря пристрастей теж описана як рух усередині людини, а не як те, що приходить ззовні. Наше порівняння — за образом, не за джерелом.}
\vfig{ilustratsii/8-1-knot.pdf}{Вузол полону: шість рухів розуму}{Схема наша, не історична ілюстрація.}
\vfig{ilustratsii-naratyvni/8-1-polon.pdf}{Людина за ґратами тягнеться рукою до птаха, що сидить на волі}{Наша ілюстрація до вірша 8.1 — не історичне зображення.}

\vsec{8.2}
\vdeva{तदा मुक्तिर् यदा चित्तं न वाञ्छति न शोचति । न मुञ्चति न गृह्णाति न हृष्यति न कुप्यति ॥}
\viast{tadā muktir yadā cittaṃ na vāñchati na śocati | na muñcati na gṛhṇāti na hṛṣyati na kupyati ||}
\vtranslit{тада муктір яда чіттам на ваньччаті на шочаті | на муньчаті на ґріхнаті на хрішяті на куп'яті ||}
\vperek{Визволення — тоді, коли розум не бажає, не сумує, не відпускає й не бере, не радіє й не гнівається.}
\vcom{Построкове зіставлення: редакторська чернетка, не рецензовано.}
\vcom{Рядок а}
\vdeva{तदा मुक्तिर् यदा चित्तं न वाञ्छति न शोचति}
\viast{tadā muktir yadā cittaṃ na vāñchati na śocati}
\vtranslit{тада муктір яда чіттам на ваньччаті на шочаті}
\vperek{Визволення — тоді, коли розум не бажає, не сумує,}
\vcom{Кореляція tadā ... yadā пов'язує визволення зі станом розуму. Форми na vāñchati й na śocati називають перші два заперечені рухи — «не бажає» і «не сумує». Це початок одного речення, яке продовжується в другому піврядку.}
\vcom{Рядок б}
\vdeva{न मुञ्चति न गृह्णाति न हृष्यति न कुप्यति}
\viast{na muñcati na gṛhṇāti na hṛṣyati na kupyati}
\vtranslit{на муньчаті на ґріхнаті на хрішяті на куп'яті}
\vperek{не відпускає й не бере, не радіє й не гнівається.}
\vcom{Чотири повтори na завершують симетричний перелік: не відпускає, не бере, не радіє, не гнівається. Порівняння з 8.1 є редакторським зіставленням двох чинних карток; 8.2 має hṛṣyati, що не дає підстав автоматично замінювати dhṛṣyati у 8.1.}
\vcom{Вірш є **дослівним запереченням** попереднього: ті самі шість дієслів, кожне з часткою na. Це важливо для перекладу: будь-яка стилістична варіація тут зруйнує паралелізм, який і є аргументом. Визволення не додає нічого — воно знімає рівно те, що було названо полоном.}
\vpar{Українська паралель (наша інтерпретація)}{Леся Українка, «Contra spem spero!» (1890).}{Наша паралель: у Лесі Українки рівновага теж не залежить від того, чи є надія. Наше порівняння — типологічне: обидва тексти описують стан, а не здобуток.}
\vfig{ilustratsii/8-2-knot.pdf}{Той самий вузол, але кожен рух заперечено}{Схема наша, не історична ілюстрація.}

\vsec{8.3}
\vdeva{तदा बन्धो यदा चित्तं सक्तं कास्वपि दृष्टिषु । तदा मोक्षो यदा चित्तम् असक्तं सर्वदृष्टिषु ॥}
\viast{tadā bandho yadā cittaṃ saktaṃ kāsvapi dṛṣṭiṣu | tadā mokṣo yadā cittam asaktaṃ sarvadṛṣṭiṣu ||}
\vtranslit{тада бандго яда чіттам сактам касвапі дріштішу | тада мокшо яда чіттам асактам сарвадріштішу ||}
\vperek{Полон — тоді, коли розум прив'язаний до якихось поглядів; визволення — коли розум неприв'язаний до всіх поглядів.}
\vcom{Построкове зіставлення: редакторська чернетка, не рецензовано.}
\vcom{Рядок а}
\vdeva{तदा बन्धो यदा चित्तं सक्तं कास्वपि दृष्टिषु}
\viast{tadā bandho yadā cittaṃ saktaṃ kāsvapi dṛṣṭiṣu}
\vtranslit{тада бандго яда чіттам сактам касвапі дріштішу}
\vperek{Полон — тоді, коли розум прив'язаний до якихось поглядів;}
\vcom{Фраза tadā bandho yadā формує умову полону; saktaṃ означає прив'язаність розуму, kāsvapi dṛṣṭiṣu — «до якихось поглядів». Це перша половина протиставлення, а не твердження про історичний вплив на українських авторів.}
\vcom{Рядок б}
\vdeva{तदा मोक्षो यदा चित्तम् असक्तं सर्वदृष्टिषु}
\viast{tadā mokṣo yadā cittam asaktaṃ sarvadṛṣṭiṣu}
\vtranslit{тада мокшо яда чіттам асактам сарвадріштішу}
\vperek{визволення — коли розум неприв'язаний до всіх поглядів.}
\vcom{Друга половина паралельна першій: tadā mokṣo yadā вводить визволення, asaktaṃ — неприв'язаність, sarvadṛṣṭiṣu — «до всіх поглядів». Контраст kāsvapi / sarva передає різний обсяг прив'язаності; переклад збережено дослівно.}
\vcom{`dṛṣṭi` — «погляд» у подвійному сенсі: і зір, і переконання. Асиметрія навмисна: `kāsvapi` («до якихось») проти `sarva` («до всіх»). Тобто достатньо **однієї** прив'язаності, щоб бути в полоні, але потрібна неприв'язаність до **всіх** — і це не той самий обсяг вимог, що звучить суворіше, ніж здається з першого читання.}
\vpar{Українська паралель (наша інтерпретація)}{Микола Хвильовий, «Я (Романтика)» (1924).}{Наша паралель: у Хвильового привʼязаність до власного погляду теж постає як те, що тримає міцніше за будь-яку зовнішню силу. Наше порівняння — за будовою думки.}
\vfig{ilustratsii/8-3-balance.pdf}{Дві шальки: привʼязаність до поглядів і свобода від них}{Схема наша, не історична ілюстрація.}

\vsec{8.4}
\vdeva{यदा नाहं तदा मोक्षो यदाहं बन्धनं तदा । मत्वेति हेलया किञ्चिन् मा गृहाण विमुञ्च मा ॥}
\viast{yadā nāhaṃ tadā mokṣo yadāhaṃ bandhanaṃ tadā | matveti helayā kiñcin mā gṛhāṇa vimuñca mā ||}
\vtranslit{яда нахам тада мокшо ядахам бандганам тада | матветі хелая кіньчін ма ґріхана вімуньча ма ||}
\vperek{Коли немає «я» — визволення; коли є «я» — полон. Так розуміючи, легко нічого не бери й нічого не відпускай.}
\vcom{Построкове зіставлення: редакторська чернетка, не рецензовано.}
\vcom{Рядок а}
\vdeva{यदा नाहं तदा मोक्षो यदाहं बन्धनं तदा}
\viast{yadā nāhaṃ tadā mokṣo yadāhaṃ bandhanaṃ tadā}
\vtranslit{яда нахам тада мокшо ядахам бандганам тада}
\vperek{Коли немає «я» — визволення; коли є «я» — полон.}
\vcom{Протиставлення yadā nāhaṃ / yadāhaṃ пов'язує відсутність чи присутність «я» з визволенням та полоном. Це філософське твердження тексту, а не психологічний діагноз; чинний український відповідник лишається без змін.}
\vcom{Рядок б}
\vdeva{मत्वेति हेलया किञ्चिन् मा गृहाण विमुञ्च मा}
\viast{matveti helayā kiñcin mā gṛhāṇa vimuñca mā}
\vtranslit{матветі хелая кіньчін ма ґріхана вімуньча ма}
\vperek{Так розуміючи, легко нічого не бери й нічого не відпускай.}
\vcom{matveti вводить розуміння попередньої тези, helayā — відтінок легкості. Подвійне mā при наказових формах gṛhāṇa та vimuñca передає дві заборони: не брати й не відпускати; жодних нових санскритських читань тут не встановлюємо.}
\vcom{Тут уперше названо причину, а не симптом: попередні три вірші описували рухи розуму, цей називає їхнє джерело — `aham`. `helayā` («легко, жартома») різко змінює тон: те, що виглядало як велика вимога, подане як щось майже незначне. І `mā … mā` — подвійна заборона: не бери й не відпускай, тобто не роби й самого жесту вибору.}
\vpar{Українська паралель (наша інтерпретація)}{Іван Багряний, «Сад Гетсиманський» (1950).}{Наша паралель: у Багряного теж є думка, що остання вʼязниця людини — вона сама. Наше порівняння — типологічне, не твердження про вплив.}
\vfig{ilustratsii/8-4-neti.pdf}{Я і його відсутність: легко не брати й не відпускати}{Схема наша, не історична ілюстрація.}

